% Section: Related Work

\paragraph*{Dynamic stopping in BCI}
Prior work uses hand-tuned thresholds on a confidence statistic,
e.g.\ SPRT-style stopping for motor
imagery~\cite{Liu2017MotorImageryBTSPRT}. None learn a policy or
include a recalibration or abstention action.

\paragraph*{Sequential decision frameworks for BCI}
\textsc{EEG\_RL-Net}~\cite{Aung2024EEG_RL_Net} trains a Dueling
DQN classification head online; error-potential (ErrP)-driven online
RL~\cite{Fidencio2025ErrPRL} adapts a decoder during a paced game.
Neither is \emph{offline} RL, neither exposes \textsc{Defer},
\textsc{Recal}, or \textsc{Abstain}, and neither ships a calibrated
OPE estimator.

\paragraph*{Offline RL and OPE methodology}
Conservative Q-Learning~\cite{Kumar2020CQL} and risk-sensitive
distributional variants~\cite{Dabney2018QR-DQN} are now standard;
PDIS~\cite{Precup2000PDIS}, doubly-robust
estimators~\cite{Jiang2016DR}, and fitted-Q
evaluation~\cite{Le2019BatchRL} are the workhorse OPE methods.
Our contribution is porting this stack to BCI \emph{and} validating
it against on-policy Monte-Carlo ground truth on a deterministic
per-trial environment.

\paragraph*{EEG foundation models}
\labram{}~\cite{LaBraM2024}, CBraMod, and EEGPT release pretrained
encoders; every published downstream usage attaches a supervised
classification head. Our cross-subject experiments
(\S\ref{sec:results:labram_loso}) are the first end-to-end RL
fine-tuning of an EEG foundation model in BCI.
